%%%%%%%%%%%%%%%%%%%%%%%%%%%%%%%%%%%%%%%%%%%%%
\section{Conclusions}
\label{Sec:Conclusions}
%%%%%%%%%%%%%%%%%%%%%%%%%%%%%%%%%%%%%%%%%%%%%

In these notes we have studied numerically the relaxation through phase space mixing of a spherically symmetric collisionless gas which is trapped in an external central potential and whose initial DF is concentrated about one orbital phase, taking into account the self-gravity of the gas, in the reduced model in which all the gas particles have the same magnitude of the angular momentum. Our analysis is based on a PIC code whose particles are placed on a regular lattice in action-angle variables. In this section we summarize the results and list some open problems.

In the absence of self-gravity, the PIC code reproduces the exact solution of the mixing problem to within $6\times 10^{-9}$ in relative terms over $10^4$ time units, during which the first harmonic of the DF decays by seven orders of magnitude. When the self-gravity of the gas is taken into account, the observables which measure the mixing have to be computed with the action-angle variables of the total potential, or of the instantaneous potential when the gas is far from equilibrium. The action-angle variables of the external potential leave a time-independent residue which is proportional to the mass of the gas ($2\times 10^{-3}$ for $M_{gas}/M = 10^{-3}$) and which could be mistaken for an incomplete mixing.

The self-gravity of the cloud changes its relaxation in three ways, which are controlled by the parameter~(\ref{Eq:muDef}) that compares the amplitude of the potential generated by the cloud with the spread of its orbital frequencies. First, it slows down the relaxation: the potential which the gas generates approaches a stationary profile while the DF continues to mix, which is the analogue of Landau damping, but the time in which its deviation falls below $10^{-3}$ of its largest value grows from $2.9$ units of the mixing time $\tau_1$ without self-gravity to $6.2\,\tau_1$ for $\mu = 0.43$. Second, it changes the state which is reached. In the absence of self-gravity this state is the average of the initial DF over the angle variables. With self-gravity the action variables of the gas particles are not conserved while the potential of the cloud oscillates, and gas accumulates on nearly circular orbits: for $\mu = 0.31$, which corresponds to a cloud with one percent of the central mass, the fraction of the mass with $J < 0.2\sigma_J$ is $6.1$ percent instead of $0.5$ percent, although the gas is completely mixed. For $\mu\gtrsim 0.4$ a compact part of the cloud is transported to the circular orbit before it is mixed and survives there as a dense core, in which the DF exceeds the largest value that can be obtained by mixing at fixed action; for $\mu = 0.43$ the core contains $10$ percent of the mass of the gas, and its radial oscillation is damped with a decay time of about $4\,\tau_1$. Third, above a critical value $\mu_c$ the oscillation of the core is not damped during the $30$ units of $\tau_1$ covered by our simulations, and the potential does not relax. For clouds which are compact in angle, $\mu_c$ lies between $0.43$ and $0.48$ in the isochrone potential as well as in the Kepler potential, and the simulations with other widths in action are compatible with this interval, whereas for a cloud which is wide in angle the transition lies between $\mu = 0.52$ and $\mu = 0.69$. In the Kepler potential this means that, for a cloud of a given shape, the relaxation by phase mixing described in Ref.~\cite{sCjL26} only takes place below a critical mass, which lies between $0.02$ and $0.04$ times the central mass for the clouds considered here.

Several questions remain open. A theory of the critical value $\mu_c$ is missing, which should also explain why a cloud which is wide in angle requires a larger value; the heuristic argument of section~\ref{SubSec:FarRuns} only shows that $\mu_c$ should be of order one. The mechanism which damps the oscillation of the core below $\mu_c$ and sustains it above, and its relation to the trapping of resonant particles in plasma waves, remains to be analyzed. Except for the largest masses, our simulations do not show whether the core is held together by its own gravity. The dependence of the results on the grid spacing, on the time step and on the number of rows of the lattice close to the transition has not been tested, and the decay laws of $|h_1|$ and of the potential require simulations with more particles than those used here. It is not known whether the oscillation of the core above $\mu_c$ is eventually damped, as in the core-halo states of systems with long-range interactions~\cite{yLrPfRtTfB14}, and on which time scale; answering this question requires longer simulations. Finally, all our results refer to the reduced model, and the core forms at the circular orbit, which is excluded from the support of the data considered in Ref.~\cite{sCjL26}. When the angular momenta of the gas particles are distributed, the frequencies of the orbits depend also on $L$, which provides an additional source of mixing (see section~\ref{Sec:Linear}), and the initial data can be chosen to satisfy the assumptions of that reference. Whether a core forms in that case, and whether the critical value of $\mu$ is shifted, is a natural next step, for which the methods of these notes can be extended in a direct way. \EJ{To be completed after the discussion with Olivier: the outlook regarding the distribution of angular momenta, and the reference to the companion article on perturbations of steady states.}
